% Ad Atticum 6.8 (ad-atticum-06-08)
% Source: https://raw.githubusercontent.com/PerseusDL/canonical-latinLit/master/data/phi0474/phi057/phi0474.phi057.perseus-lat1.xml
% Fetched by scripts/fetch_latin.py

% Dateline (Perseus): Scr. Ephes i K. Oct. a. 704 (50).

\ciceroLetterOpener{CICERO ATTICO salutem}

\ciceroSection{1} cum instituissem ad te scribere calamumque sumpsissem, Batonius e navi recta ad me venit domum Ephesi et epistulam tuam reddidit pridie Kal. Octobris. laetatus sum felicitate navigationis tuae, opportunitate Piliae, etiam hercule sermone eiusdem de coniugio Tulliae meae.

\ciceroSection{2} Batonius autem meros terrores ad me attulit Caesarianos, cum Lepta etiam plura locutus est, spero falsa, sed certe horribilia, exercitum nullo modo dimissurum, cum illo praetores designatos, Cassium tribunum pl., Lentulum consulem facere, Pompeio in animo esse urbem relinquere.

\ciceroSection{3} sed heus tu! numquid moleste fers de illo qui se solet anteferre patruo sororis tuae fili? at a quibus victus! sed ad rem.

\ciceroSection{4} nos etesiae vehementissime tardarunt; detraxit xx ipsos dies etiam aphractus Rhodiorum. Kal. Octobr. Epheso conscendentes hanc epistulam dedimus L. Tarquitio simul e portu egredienti sed expeditius naviganti. nos Rhodiorum aphractis ceterisque longis navibus tranquillitates aucupaturi eramus; ita tamen properabamus ut non posset magis.

\ciceroSection{5} de raudusculo Puteolano gratum. nunc velim dispicias res Romanas, videas quid nobis de triumpho cogitandum putes ad quem amici me vocant. ego nisi Bibulus qui, dum unus hostis in Syria fuit, pedem porta non plus extulit quam domo sua, adniteretur de triumpho, aequo animo essem. nunc vero αἰσχρὸν . sed explora rem totam, ut quo die congressi erimus consilium capere possimus. sat multa, qui et properarem et ei litteras darem qui aut mecum aut paulo ante venturus esset. Cicero tibi plurimam salutem dicit. tu dices utriusque nostrum verbis et Piliae tuae et filiae.
